% GENERATED FILE. Edit canonical lessons, then run npm run generate:books.
% canonical-source-hash: fnv1a64:b8ea5399b31fd515
% canonical-lessons: GU-C424-kesiyar, GU-C424-pasbuk, GU-C424-kapli, GU-C424-tokan, GU-C424-kard, GU-R424-review-417-424

\chapter{At the Bank Counter}
\label{ch:gu-at-the-bank-counter}
\hlchaptermodality{424}

\begin{chapteropening}
\textbf{By the end of this chapter:} \emph{I can do simple business at a bank counter.}

Complete the last lesson of chapter 424: I can do simple business at a bank counter.
\end{chapteropening}

\section[\texorpdfstring{\gu{કેશિયર} (keśiyar)}{કેશિયર (keśiyar)}]{\gu{કેશિયર} (keśiyar) — a cashier}

\label{lesson:GU-C424-kesiyar}



\begin{quote}
Before the new one: say the Gujarati for a stink, then the Gujarati for to set up a home.
\end{quote}

\subsection*{You'll want to know: \gu{કેશિયર}}
\textbf{\gu{કેશિયર}} — \emph{keśiyar} — \textquotedblleft{}a cashier\textquotedblright{}.

\begin{quote}
\gu{બેંકના} \gu{કેશિયરે} \gu{મને} \gu{નવી} \gu{નોટો} \gu{આપી}.
\end{quote}

\begin{quote}
\emph{beṅknā keśiyre mane navī noṭo āpī.}
\end{quote}

\begin{quote}
The bank cashier gave me new notes.
\end{quote}

A masculine noun.

\begin{cousinweb}
\textbf{\gu{કેશિયર}} (\emph{keśiyar}) is borrowed from English \emph{cashier}.
\end{cousinweb}

\begin{grammarlens}[title={What you've built}]
One more word for getting around.
\end{grammarlens}

\subsection*{Your turn}
\begin{itemize}
\raggedright
  \item \emph{Say it:} \emph{keśiyar}
  \item \emph{Say it:} \emph{keśiyar}, once more
  \item \emph{Say it:} \emph{beṅknā keśiyre mane navī noṭo āpī.} — The bank cashier gave me new notes.
  \item \emph{Recall:} say the Gujarati for a stink, then the Gujarati for to set up a home, then say \emph{keśiyar} again
\end{itemize}

\begin{tcolorbox}[breakable,colback=teal!4,colframe=teal!35!black,title={Before you move on}]
What does \gu{કેશિયર} mean? (\textquotedblleft{}A cashier\textquotedblright{}.) Say it once more.
\end{tcolorbox}
\section[\texorpdfstring{\gu{પાસબુક} (pāsbuk)}{પાસબુક (pāsbuk)}]{\gu{પાસબુક} (pāsbuk) — a passbook (of a bank account)}

\label{lesson:GU-C424-pasbuk}



\begin{quote}
Before the new one: say the Gujarati for to set up a home, then the Gujarati for a cashier.
\end{quote}

\subsection*{You'll want to know: \gu{પાસબુક}}
\textbf{\gu{પાસબુક}} — \emph{pāsbuk} — \textquotedblleft{}a passbook (of a bank account)\textquotedblright{}.

\begin{quote}
\gu{મારી} \gu{પાસબુક} \gu{ઘરે} \gu{રહી} \gu{ગઈ}.
\end{quote}

\begin{quote}
\emph{mārī pāsbuk ghare rahī gaī.}
\end{quote}

\begin{quote}
I left my passbook at home.
\end{quote}

A feminine noun.

\begin{cousinweb}
\textbf{\gu{પાસબુક}} (\emph{pāsbuk}) is borrowed from English \emph{passbook}, the little book where a bank writes up your account.
\end{cousinweb}

\begin{grammarlens}[title={What you've built}]
One more word for getting around.
\end{grammarlens}

\subsection*{Your turn}
\begin{itemize}
\raggedright
  \item \emph{Say it:} \emph{pāsbuk}
  \item \emph{Say it:} \emph{pāsbuk}, once more
  \item \emph{Say it:} \emph{mārī pāsbuk ghare rahī gaī.} — I left my passbook at home.
  \item \emph{Recall:} say the Gujarati for to set up a home, then the Gujarati for a cashier, then say \emph{pāsbuk} again
\end{itemize}

\begin{tcolorbox}[breakable,colback=teal!4,colframe=teal!35!black,title={Before you move on}]
What does \gu{પાસબુક} mean? (\textquotedblleft{}A passbook (of a bank account)\textquotedblright{}.) Say it once more.
\end{tcolorbox}
\section[\texorpdfstring{\gu{કાપલી} (kāplī)}{કાપલી (kāplī)}]{\gu{કાપલી} (kāplī) — a slip of paper}

\label{lesson:GU-C424-kapli}



\begin{quote}
Before the new one: say the Gujarati for a cashier, then the Gujarati for a passbook.
\end{quote}

\subsection*{You'll want to know: \gu{કાપલી}}
\textbf{\gu{કાપલી}} — \emph{kāplī} — \textquotedblleft{}a slip of paper\textquotedblright{}.

\begin{quote}
\gu{પૈસા} \gu{જમા} \gu{કરવા} \gu{માટે} \gu{આ} \gu{કાપલી} \gu{ભરો}.
\end{quote}

\begin{quote}
\emph{paisā jamā karvā māṭe ā kāplī bharo.}
\end{quote}

\begin{quote}
Fill in this slip to deposit money.
\end{quote}

A feminine noun.

\begin{cousinweb}
\textbf{\gu{કાપલી}} (\emph{kāplī}) is built on \gu{કાપવું}, to cut, a tadbhava of Sanskrit \emph{kalpayati}, \textquotedblleft{}trims, cuts\textquotedblright{} (CDIAL 2944): a strip of paper cut off.
\end{cousinweb}

\begin{grammarlens}[title={What you've built}]
One more word for getting around.
\end{grammarlens}

\subsection*{Your turn}
\begin{itemize}
\raggedright
  \item \emph{Say it:} \emph{kāplī}
  \item \emph{Say it:} \emph{kāplī}, once more
  \item \emph{Say it:} \emph{paisā jamā karvā māṭe ā kāplī bharo.} — Fill in this slip to deposit money.
  \item \emph{Recall:} say the Gujarati for a cashier, then the Gujarati for a passbook, then say \emph{kāplī} again
\end{itemize}

\begin{tcolorbox}[breakable,colback=teal!4,colframe=teal!35!black,title={Before you move on}]
What does \gu{કાપલી} mean? (\textquotedblleft{}A slip of paper\textquotedblright{}.) Say it once more.
\end{tcolorbox}
\section[\texorpdfstring{\gu{ટોકન} (ṭokan)}{ટોકન (ṭokan)}]{\gu{ટોકન} (ṭokan) — a token (a numbered slip for your turn)}

\label{lesson:GU-C424-tokan}



\begin{quote}
Before the new one: say the Gujarati for a passbook, then the Gujarati for a slip of paper.
\end{quote}

\subsection*{You'll want to know: \gu{ટોકન}}
\textbf{\gu{ટોકન}} — \emph{ṭokan} — \textquotedblleft{}a token (a numbered slip for your turn)\textquotedblright{}.

\begin{quote}
\gu{બેંકમાં} \gu{મને} \gu{પાંચ} \gu{નંબરનું} \gu{ટોકન} \gu{મળ્યું}.
\end{quote}

\begin{quote}
\emph{beṅkmā̃ mane pā̃c nambarnũ ṭokan maḷyũ.}
\end{quote}

\begin{quote}
At the bank I got token number five.
\end{quote}

A neuter noun.

\begin{cousinweb}
\textbf{\gu{ટોકન}} (\emph{ṭokan}) is borrowed from English \emph{token}: at a bank or a clinic, the numbered slip that keeps your place in the queue.
\end{cousinweb}

\begin{grammarlens}[title={What you've built}]
One more word for getting around.
\end{grammarlens}

\subsection*{Your turn}
\begin{itemize}
\raggedright
  \item \emph{Say it:} \emph{ṭokan}
  \item \emph{Say it:} \emph{ṭokan}, once more
  \item \emph{Say it:} \emph{beṅkmā̃ mane pā̃c nambarnũ ṭokan maḷyũ.} — At the bank I got token number five.
  \item \emph{Recall:} say the Gujarati for a passbook, then the Gujarati for a slip of paper, then say \emph{ṭokan} again
\end{itemize}

\begin{tcolorbox}[breakable,colback=teal!4,colframe=teal!35!black,title={Before you move on}]
What does \gu{ટોકન} mean? (\textquotedblleft{}A token (a numbered slip for your turn)\textquotedblright{}.) Say it once more.
\end{tcolorbox}
\section[\texorpdfstring{\gu{કાર્ડ} (kārḍ)}{કાર્ડ (kārḍ)}]{\gu{કાર્ડ} (kārḍ) — a card (a bank card)}

\label{lesson:GU-C424-kard}



\begin{quote}
Before the new one: say the Gujarati for a slip of paper, then the Gujarati for a token.
\end{quote}

\subsection*{You'll want to know: \gu{કાર્ડ}}
\textbf{\gu{કાર્ડ}} — \emph{kārḍ} — \textquotedblleft{}a card (a bank card)\textquotedblright{}.

\begin{quote}
\gu{તેણે} \gu{કાર્ડથી} \gu{બિલ} \gu{ચૂકવ્યું}.
\end{quote}

\begin{quote}
\emph{teṇe kārḍathī bil cūkavyũ.}
\end{quote}

\begin{quote}
She paid the bill by card.
\end{quote}

A neuter noun.

\begin{cousinweb}
\textbf{\gu{કાર્ડ}} (\emph{kārḍ}) is borrowed from English \emph{card}, here a bank card.
\end{cousinweb}

\begin{grammarlens}[title={What you've built}]
One more word for getting around.
\end{grammarlens}

\subsection*{Your turn}
\begin{itemize}
\raggedright
  \item \emph{Say it:} \emph{kārḍ}
  \item \emph{Say it:} \emph{kārḍ}, once more
  \item \emph{Say it:} \emph{teṇe kārḍathī bil cūkavyũ.} — She paid the bill by card.
  \item \emph{Recall:} say the Gujarati for a slip of paper, then the Gujarati for a token, then say \emph{kārḍ} again
\end{itemize}

\begin{tcolorbox}[breakable,colback=teal!4,colframe=teal!35!black,title={Before you move on}]
What does \gu{કાર્ડ} mean? (\textquotedblleft{}A card (a bank card)\textquotedblright{}.) Say it once more.
\end{tcolorbox}
\section[(dialogue)]{Looking back: from The Pol to At the Bank Counter}

\label{lesson:GU-R424-review-417-424}



\begin{quote}
Say the Gujarati for a pol and an entrance gateway.
\end{quote}

\subsection*{Your turn}
Say these aloud:

\begin{itemize}
\raggedright
  \item a pol, an entrance gateway, a raised platform, a small gated lane, to shut (a door)
  \item a housing society, a flat, a bungalow, a water tank, a tenant
  \item a street, a neighbourhood lane, a quarter of a town, a resident, to wander aimlessly
  \item a wooden door, a roof tile, a hedge, a post, to thatch
  \item a din, an uproar, cattle, to growl, to chase away
  \item a handcart, a hawker, a paan stall, a sweet-maker, to sell well
  \item a town, a district, a sewer, a stink, to set up a home
  \item a cashier, a passbook, a slip of paper, a token, a bank card
\end{itemize}

\begin{tcolorbox}[breakable,colback=teal!4,colframe=teal!35!black,title={Before you move on}]
A pol? (\textbf{\gu{પોળ}}, \emph{poḷ}.) An entrance gateway? (\textbf{\gu{ડેલી}}, \emph{ḍelī}.) A raised platform? (\textbf{\gu{ઓટલો}}, \emph{oṭlo}.) A small gated lane? (\textbf{\gu{ખડકી}}, \emph{khaḍkī}.) To shut (a door)? (\textbf{\gu{વાસવું}}, \emph{vāsvũ}.) A housing society? (\textbf{\gu{સોસાયટી}}, \emph{sosāyṭī}.) A flat? (\textbf{\gu{ફ્લેટ}}, \emph{phleṭ}.) A bungalow? (\textbf{\gu{બંગલો}}, \emph{banglo}.) A water tank? (\textbf{\gu{ટાંકી}}, \emph{ṭā̃kī}.) A tenant? (\textbf{\gu{ભાડૂત}}, \emph{bhāḍūt}.) A street? (\textbf{\gu{શેરી}}, \emph{śerī}.) A neighbourhood lane? (\textbf{\gu{ફળિયું}}, \emph{phaḷiyũ}.) A quarter of a town? (\textbf{\gu{મહોલ્લો}}, \emph{mahollo}.) A resident? (\textbf{\gu{રહેવાસી}}, \emph{rahevāsī}.) To wander aimlessly? (\textbf{\gu{રઝળવું}}, \emph{rajhaḷvũ}.) A wooden door? (\textbf{\gu{કમાડ}}, \emph{kamāḍ}.) A roof tile? (\textbf{\gu{નળિયું}}, \emph{naḷiyũ}.) A hedge? (\textbf{\gu{વાડ}}, \emph{vāḍ}.) A post? (\textbf{\gu{થાંભલો}}, \emph{thā̃bhlo}.) To thatch? (\textbf{\gu{છાવું}}, \emph{chāvũ}.) A din? (\textbf{\gu{ઘોંઘાટ}}, \emph{ghõghāṭ}.) An uproar? (\textbf{\gu{હોબાળો}}, \emph{hobāḷo}.) Cattle? (\textbf{\gu{ઢોર}}, \emph{ḍhor}.) To growl? (\textbf{\gu{ઘુરકવું}}, \emph{ghurakvũ}.) To chase away? (\textbf{\gu{ભગાડવું}}, \emph{bhagāḍvũ}.) A handcart? (\textbf{\gu{લારી}}, \emph{lārī}.) A hawker? (\textbf{\gu{ફેરિયો}}, \emph{pheriyo}.) A paan stall? (\textbf{\gu{ગલ્લો}}, \emph{gallo}.) A sweet-maker? (\textbf{\gu{કંદોઈ}}, \emph{kandoī}.) To sell well? (\textbf{\gu{ખપવું}}, \emph{khapvũ}.) A town? (\textbf{\gu{નગર}}, \emph{nagar}.) A district? (\textbf{\gu{જિલ્લો}}, \emph{jillo}.) A sewer? (\textbf{\gu{ગટર}}, \emph{gaṭar}.) A stink? (\textbf{\gu{દુર્ગંધ}}, \emph{durgandh}.) To set up a home? (\textbf{\gu{વસાવવું}}, \emph{vasāvvũ}.) A cashier? (\textbf{\gu{કેશિયર}}, \emph{keśiyar}.) A passbook? (\textbf{\gu{પાસબુક}}, \emph{pāsbuk}.) A slip of paper? (\textbf{\gu{કાપલી}}, \emph{kāplī}.) A token? (\textbf{\gu{ટોકન}}, \emph{ṭokan}.) A bank card? (\textbf{\gu{કાર્ડ}}, \emph{kārḍ}.)
\end{tcolorbox}
